\documentclass[11pt]{article}

\usepackage[margin=1in]{geometry}
\usepackage[T1]{fontenc}
\usepackage[utf8]{inputenc}
\usepackage{lmodern}
\usepackage{microtype}
\usepackage{amsmath, amssymb, amsthm, mathtools}
\usepackage{booktabs}
\usepackage{enumitem}
\usepackage{xcolor}
\usepackage[colorlinks=true, linkcolor=blue!55!black, citecolor=blue!55!black, urlcolor=blue!55!black]{hyperref}

% -----------------------------------------------------------------------------
% Notation macros
% -----------------------------------------------------------------------------

\newcommand{\R}{\mathbb{R}}
\newcommand{\E}{\mathbb{E}}
\newcommand{\ip}[2]{\left\langle #1, #2 \right\rangle}
\newcommand{\norm}[1]{\left\lVert #1 \right\rVert}
\newcommand{\softmax}{\operatorname{softmax}}
\newcommand{\diag}{\operatorname{diag}}

\newcommand{\dm}{d_{\mathrm{model}}}
\newcommand{\dhead}{d_{\mathrm{head}}}
\newcommand{\dsae}{d_{\mathrm{SAE}}}

\newcommand{\pretag}{\mathrm{pre}}
\newcommand{\midtag}{\mathrm{mid}}
\newcommand{\attn}{\mathrm{attn}}
\newcommand{\OV}{\mathrm{OV}}
\newcommand{\QK}{\mathrm{QK}}
\newcommand{\LN}{\mathrm{LN}}

\newcommand{\hook}[1]{\texttt{#1}}
\newcommand{\residpre}{r^{\pretag}}
\newcommand{\residmid}{r^{\midtag}}
\newcommand{\attnout}{\mathrm{attn\_out}}

\newcommand{\todo}[1]{\textcolor{red!65!black}{\textbf{TODO:} #1}}

% -----------------------------------------------------------------------------
% Theorem-like environments
% -----------------------------------------------------------------------------

\theoremstyle{definition}
\newtheorem{definition}{Definition}
\newtheorem{protocol}{Protocol}

\theoremstyle{plain}
\newtheorem{proposition}{Proposition}
\newtheorem{claim}{Claim}

\theoremstyle{remark}
\newtheorem{remark}{Remark}

% -----------------------------------------------------------------------------

\title{Feature-Resolved Attention: A Clean Derivation}
\author{Dmitry Manning-Coe}
\date{\today}

\begin{document}

\maketitle

\begin{abstract}
Feature-resolved attention (FRA) decomposes an attention circuit in the basis
of sparse autoencoder (SAE) features.  The goal is to assign a scalar target,
such as a downstream feature preactivation or steering direction, to specific
query, key, and value-side feature contributions.  This note gives a clean
derivation of the two pieces we use most often: frozen-attention OV attribution
and first-order QK attribution via the softmax Jacobian.  It also records the
distinction between feature ranking and intervention path, since conflating
these two choices has been a recurring source of confusion.
\end{abstract}

\tableofcontents

% =============================================================================
\section{Roadmap}
% =============================================================================

The derivation separates three questions.

\begin{enumerate}[leftmargin=*]
  \item \textbf{What scalar are we explaining?}
  We choose a target direction \(d \in \R^{\dm}\) and explain
  \[
    T_q(d) := \ip{\attnout_q}{d}.
  \]
  In the sleeper setting, \(d\) can be the encoder direction of a downstream
  suppressor feature at \hook{hook\_resid\_mid}.

  \item \textbf{How do features contribute through a fixed attention pattern?}
  Holding \(A\) fixed makes the OV path linear.  This gives a clean
  per-head, per-source-token, per-feature attribution.

  \item \textbf{How do features change the attention pattern?}
  Attention patterns are nonlinear in QK scores.  We therefore decompose the
  scores exactly in feature pairs and then use a first-order softmax-Jacobian
  approximation to attribute changes in the scalar target.
\end{enumerate}

\paragraph{Main caution.}
An FRA \emph{ranking} and an FRA \emph{intervention path} are separate design
choices.  For example, one can rank features by QK attribution but intervene
through the V path, the QK path, or the original normalized residual hook.

% =============================================================================
\section{Transformer and Hook Conventions}
% =============================================================================

We use the TransformerLens row-vector convention: activations multiply weights
on the right.  Thus an activation update has the form
\[
  \texttt{new\_activation} = \texttt{old\_activation}\, W + b.
\]

Fix a layer \(\ell\), head \(h\), query position \(q\), and source/key
position \(k\).  Let
\[
  x_t \in \R^{\dm}
\]
denote the exact tensor read by the attention sublayer at position \(t\).  In
TransformerLens terms this is typically the layer-normalized attention input,
for example \hook{blocks.\(\ell\).ln1.hook\_normalized}, after whatever
folding convention the model uses.

For head \(h\), define
\[
  Q_t^h = x_t W_Q^h,\qquad
  K_t^h = x_t W_K^h,\qquad
  V_t^h = x_t W_V^h + b_V^h.
\]

The pre-softmax score and attention pattern are
\[
  s_{qk}^h =
  \frac{Q_q^h \cdot K_k^h}{\sqrt{\dhead}},
  \qquad
  A_{qk}^h = \softmax_k(s_{qk}^h).
\]

The head output written to the residual stream is
\[
  y_q^h
  =
  \sum_k A_{qk}^h V_k^h W_O^h
  =
  \sum_k A_{qk}^h x_k W_V^h W_O^h + c^h,
  \qquad
  c^h := b_V^h W_O^h.
\]

The full attention write is
\[
  \attnout_q = \sum_h y_q^h + b_O.
\]

If we are targeting a direction at \hook{hook\_resid\_mid}, then
\[
  \residmid_{\ell,q}
  =
  \residpre_{\ell,q}
  +
  \attnout_{\ell,q},
\]
so the downstream feature preactivation decomposes into a skip term plus an
attention term:
\[
  \ip{\residmid_{\ell,q}}{d}
  =
  \ip{\residpre_{\ell,q}}{d}
  +
  \ip{\attnout_{\ell,q}}{d}.
\]
The FRA derivation below focuses on the attention term.

% =============================================================================
\section{Feature Decomposition}
% =============================================================================

Assume an SAE or dictionary on the attention input \(x_t\):
\[
  x_t
  =
  \sum_{\lambda=1}^{\dsae} u_t^\lambda f_\lambda
  +
  e_t.
\]
Here
\begin{itemize}[leftmargin=*]
  \item \(u_t^\lambda\) is feature \(\lambda\)'s activation at position \(t\),
  \item \(f_\lambda \in \R^{\dm}\) is the decoder direction,
  \item \(e_t\) is reconstruction error.
\end{itemize}

We will often drop \(e_t\) in the algebra and then treat reconstruction error
as an empirical error term.

\begin{remark}
The same feature basis can be used for Q, K, and V because they read the same
attention-input tensor \(x_t\).  But the channel-specific vectors
\[
  f_\lambda W_Q^h,\qquad f_\lambda W_K^h,\qquad f_\lambda W_V^h W_O^h
\]
are different.  FRA is useful precisely because it keeps these channels
separate.
\end{remark}

% =============================================================================
\section{Frozen-Attention OV Attribution}
% =============================================================================

Start with the scalar target
\[
  T_q(d) := \ip{\attnout_q}{d}.
\]
For one head, holding \(A^h\) fixed gives
\[
  \ip{y_q^h}{d}
  =
  \sum_k A_{qk}^h \ip{x_k W_V^h W_O^h}{d}
  +
  \ip{c^h}{d}.
\]

Define the virtual OV read direction \(u_h^{\OV}(d)\) by
\[
  \ip{x_k W_V^h W_O^h}{d}
  =
  \ip{x_k}{u_h^{\OV}(d)}.
\]
With row-vector activations, this is equivalently
\[
  u_h^{\OV}(d)
  =
  W_V^h W_O^h d^\top
\]
up to the usual row/column identification.

Substitute the feature expansion of \(x_k\):
\[
  \ip{x_k}{u_h^{\OV}(d)}
  \approx
  \sum_\lambda u_k^\lambda
  \underbrace{\ip{f_\lambda W_V^h W_O^h}{d}}_{:=\,\beta_{h,\lambda}(d)}.
\]

\begin{definition}[OV feature coefficient]
For target direction \(d\), define
\[
  \beta_{h,\lambda}(d)
  :=
  \ip{f_\lambda W_V^h W_O^h}{d}.
\]
\end{definition}

\begin{proposition}[Frozen-attention OV attribution]
Given the observed attention pattern \(A\), the feature-resolved OV
contribution of source feature \(\lambda\) at position \(k\), read by query
position \(q\) through head \(h\), is
\[
  C^{\OV}_{h,q,k,\lambda}(d)
  :=
  A_{qk}^h\, u_k^\lambda\, \beta_{h,\lambda}(d).
\]
Thus
\[
  T_q(d)
  \approx
  \sum_h \sum_k \sum_\lambda
  C^{\OV}_{h,q,k,\lambda}(d)
  +
  \mathrm{const}
  +
  \mathrm{reconstruction\ error}.
\]
\end{proposition}

\paragraph{Typical feature ranking.}
For a prompt distribution \(\mathcal{D}\), define an OV score
\[
  \mathrm{Score}_{\OV}(\lambda)
  =
  \sum_h
  \beta_{h,\lambda}(d)
  \;
  \E_{b \sim \mathcal{D},\, q \in \mathrm{prompt}(b)}
  \left[
    \sum_k A_{b,qk}^h u_{b,k}^{\lambda}
  \right].
\]
Rank features by \(|\mathrm{Score}_{\OV}(\lambda)|\), or by a variant that
subtracts the same statistic on a control distribution.

\todo{Specify which distribution is used for the final paper: deployment only,
deployment minus clean, or a prompt-position-restricted variant.}

% =============================================================================
\section{QK Score Decomposition}
% =============================================================================

OV attribution holds \(A\) fixed.  To ask how features change the attention
pattern, begin at the score level:
\[
  s_{qk}^h =
  \frac{(x_q W_Q^h) \cdot (x_k W_K^h)}{\sqrt{\dhead}}.
\]

Using the feature expansion on both query and key positions,
\[
  x_q \approx \sum_\mu u_q^\mu f_\mu,
  \qquad
  x_k \approx \sum_\nu u_k^\nu f_\nu.
\]
Then
\[
  s_{qk}^h
  \approx
  \sum_{\mu,\nu}
  u_q^\mu u_k^\nu
  \omega_{\mu\nu}^{h,\QK}
  +
  \mathrm{bias/error},
\]
where
\[
  \omega_{\mu\nu}^{h,\QK}
  :=
  \frac{(f_\mu W_Q^h) \cdot (f_\nu W_K^h)}{\sqrt{\dhead}}.
\]

\begin{definition}[QK score atom]
The score-level QK atom for query feature \(\mu\) and key feature \(\nu\) is
\[
  S_{qk}^{h,\mu\nu}
  :=
  u_q^\mu u_k^\nu
  \omega_{\mu\nu}^{h,\QK}.
\]
\end{definition}

\begin{remark}
The score decomposition is clean.  The attention-pattern decomposition is not,
because \(A_{qk}^h = \softmax_k(s_{qk}^h)\) is nonlinear.  This is why QK-side
target attribution needs either an exact intervention experiment or a
linearization.
\end{remark}

% =============================================================================
\section{QK Attribution via the Softmax Jacobian}
% =============================================================================

To attribute changes in the scalar target \(T_q(d)\) through QK, hold the OV
side fixed and ask how a score perturbation changes the attention-weighted OV
read.

Define the fixed-OV source scalar
\[
  g_k^h(d)
  :=
  \sum_\lambda u_k^\lambda \beta_{h,\lambda}(d).
\]
Then
\[
  T_q(d)
  \approx
  \sum_h \sum_k A_{qk}^h g_k^h(d)
  +
  \mathrm{const}.
\]

The softmax Jacobian is
\[
  \frac{\partial A_{qk'}^h}{\partial s_{qj}^h}
  =
  A_{qk'}^h
  \left(
    \mathbf{1}_{k'=j} - A_{qj}^h
  \right).
\]
Therefore
\[
  \frac{\partial T_q(d)}{\partial s_{qj}^h}
  =
  A_{qj}^h
  \left(
    g_j^h(d) - \bar g_q^h(d)
  \right),
  \qquad
  \bar g_q^h(d) := \sum_k A_{qk}^h g_k^h(d).
\]

\begin{definition}[Centered OV profile]
Define
\[
  \widetilde g_{qj}^h(d)
  :=
  A_{qj}^h
  \left(
    g_j^h(d) - \bar g_q^h(d)
  \right).
\]
\end{definition}

\subsection{Single query-feature sensitivity}

Remove query-side feature \(\mu\) at query position \(q\):
\[
  u_q^\mu \mapsto 0.
\]
The induced score shift at every key \(j\) is
\[
  \delta s_{qj}^h
  =
  -
  u_q^\mu
  \underbrace{
  \sum_\nu u_j^\nu \omega_{\mu\nu}^{h,\QK}
  }_{:=\,\kappa_j^{h,\mu}}.
\]
Thus the first-order target change is
\[
  \delta T_q^\mu(d)
  \approx
  -
  u_q^\mu
  \sum_h \sum_j
  \kappa_j^{h,\mu}
  \widetilde g_{qj}^h(d).
\]

\subsection{Single key-feature sensitivity}

Remove key-side feature \(\nu\) at source position \(k\):
\[
  u_k^\nu \mapsto 0.
\]
Only score \(s_{qk}^h\) changes, with
\[
  \delta s_{qk}^h
  =
  -
  u_k^\nu
  \underbrace{
  \sum_\mu u_q^\mu \omega_{\mu\nu}^{h,\QK}
  }_{:=\,\eta_q^{h,\nu}}.
\]
Therefore
\[
  \delta T_{q,k}^{\nu}(d)
  \approx
  -
  u_k^\nu
  \sum_h
  \eta_q^{h,\nu}
  \widetilde g_{qk}^h(d).
\]

\subsection{Pair sensitivity}

For the query-key feature pair \((\mu,\nu)\), the corresponding first-order
score contribution is
\[
  u_q^\mu u_k^\nu \omega_{\mu\nu}^{h,\QK}.
\]
Removing that atom gives
\[
  \delta T_{q,k}^{\mu,\nu}(d)
  \approx
  -
  u_q^\mu u_k^\nu
  \sum_h
  \omega_{\mu\nu}^{h,\QK}
  \widetilde g_{qk}^h(d).
\]

\paragraph{Typical QK feature ranking.}
For query-feature ranking, compute
\[
  \mathrm{Pred}_{\QK}[\mu](b,q)
  :=
  u_{b,q}^\mu
  \sum_h \sum_j
  \kappa_{b,j}^{h,\mu}
  \widetilde g_{b,qj}^h(d),
\]
and aggregate over prompt positions:
\[
  \mathrm{Score}_{\QK}(\mu)
  =
  \E_{(b,q)\in \mathrm{prompt}}
  \left[
    \left|\mathrm{Pred}_{\QK}[\mu](b,q)\right|
  \right].
\]
This is the ``QK L1 mean'' ranking used in the sleeper QK/QK note.

% =============================================================================
\section{Intervention Paths}
% =============================================================================

Attribution gives a ranked feature set \(F\).  Intervention path is a separate
choice.

Given selected features \(F\), define the SAE feature-removal delta at the
attention input:
\[
  \Delta x_t
  =
  \sum_{\lambda \in F}
  \left[
    \mathrm{decode}(z_t \ \mathrm{with}\ z_t^\lambda=0)
    -
    \mathrm{decode}(z_t)
  \right]
  \approx
  -
  \sum_{\lambda \in F} u_t^\lambda f_\lambda.
\]

\begin{protocol}[All-path intervention]
Add the delta directly at the attention input:
\[
  x_t \leftarrow x_t + \alpha \Delta x_t.
\]
This changes Q, K, and V simultaneously.
\end{protocol}

\begin{protocol}[OV-only intervention]
Project the delta through \(W_V^h\) and add it at the value hook:
\[
  V_t^h \leftarrow V_t^h + \alpha \Delta x_t W_V^h.
\]
This changes value vectors but leaves the attention pattern fixed.
\end{protocol}

\begin{protocol}[QK-only intervention]
Project the delta through \(W_Q^h\) and \(W_K^h\), then add it at the Q and K
hooks:
\[
  Q_t^h \leftarrow Q_t^h + \alpha \Delta x_t W_Q^h,
  \qquad
  K_t^h \leftarrow K_t^h + \alpha \Delta x_t W_K^h.
\]
This changes the attention pattern while leaving values unchanged.
\end{protocol}

\begin{remark}
This gives a useful grid: QK-ranked features can be sent through QK, OV, or all
paths; OV-ranked features can likewise be sent through QK, OV, or all paths.
The labels ``QK/QK'', ``QK/OV'', and ``OV/OV'' should therefore be read as
``ranking/intervention-path'', not as a single inseparable method.
\end{remark}

% =============================================================================
\section{Implementation Correspondence}
% =============================================================================

\begin{table}[h]
  \centering
  \begin{tabular}{lll}
    \toprule
    Math object & Typical tensor/code object & Shape sketch \\
    \midrule
    \(x_t\) & \hook{ln1.hook\_normalized} & batch, seq, \(\dm\) \\
    \(u_t^\lambda\) & SAE encodings \(z\) & batch, seq, \(\dsae\) \\
    \(f_\lambda\) & \hook{sae.W\_dec[lambda]} & \(\dsae, \dm\) \\
    \(A_{qk}^h\) & \hook{attn\_pattern} & batch, head, query, key \\
    \(W_Q^h,W_K^h,W_V^h\) & \hook{model.W\_Q[ell]}, etc. & head, \(\dm\), \(\dhead\) \\
    \(\Delta x_t\) & \hook{compute\_sae\_delta(...)} & batch, seq, \(\dm\) \\
    \bottomrule
  \end{tabular}
  \caption{Mapping between the derivation and typical TransformerLens/SAE tensors.}
\end{table}

\paragraph{Minimal QK/QK pseudocode.}
\begin{verbatim}
features = topk(qk_l1_mean_scores, k=K)
delta = sum(compute_sae_delta(feature=f) for f in features)

q_delta = einsum("btd,hdk->bthk", delta, W_Q[layer])
k_delta = einsum("btd,hdk->bthk", delta, W_K[layer])

hook_q: q[:, :P, :, :] += alpha * q_delta
hook_k: k[:, :P, :, :] += alpha * k_delta
\end{verbatim}

% =============================================================================
\section{Open Choices for the Paper}
% =============================================================================

\begin{enumerate}[leftmargin=*]
  \item \textbf{Target distribution.}
  Decide whether the headline ranking uses deployment prompts only,
  deployment-minus-clean, or a prompt-position mask.

  \item \textbf{Scalar target.}
  Decide whether to present \(d\) as a downstream SAE encoder direction, a logit
  direction, or a generic steering direction.

  \item \textbf{Error terms.}
  Decide how much of the SAE reconstruction error and LN linearization error
  should appear in the main text versus an appendix.

  \item \textbf{Intervention terminology.}
  Use consistent names:
  \[
    \text{ranking source} / \text{intervention path}
  \]
  for example QK/QK, QK/OV, OV/OV, and QK/All.
\end{enumerate}

% =============================================================================
\appendix
\section{Appendix Placeholders}
% =============================================================================

\subsection{LayerNorm linearization}
\todo{Insert the resid-pre to ln1 two-stage derivation if this writeup needs
the sleeper-specific two-stage path.}

\subsection{Channel-wise interaction decomposition}
\todo{Insert the Q/K/V coalition or Harsanyi decomposition if the paper needs
singles, pairs, and triples rather than the frozen-OV/QK split.}

\subsection{Experimental protocol}
\todo{Insert exact prompts, feature counts, alpha sweeps, and metrics for each
model family.}

\end{document}
